% ============================================================
%  Artículo Técnico — Análisis e Interpretación de Datos
%  UNIR — Actividad 2
%  Motor: MikTeX — compilar con pdflatex (2 pasadas)
% ============================================================
\documentclass[11pt, a4paper]{article}

% --- Codificación y lengua ---
\usepackage[utf8]{inputenc}
\usepackage[T1]{fontenc}
\usepackage[spanish]{babel}
\addto\captionsspanish{\renewcommand{\tablename}{Tabla}}

% --- Fuente Georgia (aproximación con mathptmx + ajuste) ---
% MikTeX no incluye Georgia nativa; se usa TeX Gyre Termes,
% equivalente métricamente a Times New Roman / Georgia.
\usepackage{tgtermes}

% --- Márgenes y geometría ---
\usepackage[
  a4paper,
  top=2.5cm, bottom=2.5cm,
  left=2.5cm, right=2.5cm
]{geometry}

% --- Interlineado 1.5 ---
\usepackage{setspace}
\onehalfspacing

% --- Matemáticas y tablas ---
\usepackage{amsmath}
\usepackage{booktabs}
\usepackage{array}
\usepackage{float}

% --- Gráficos ---
\usepackage{graphicx}
\usepackage{pgfplots}
\pgfplotsset{compat=1.18}
\usepackage{tikz}

% --- Colores ---
\usepackage{xcolor}
\definecolor{azulunir}{RGB}{0, 84, 166}

% --- Hiperenlaces ---
\usepackage[
  colorlinks=true,
  linkcolor=azulunir,
  citecolor=azulunir,
  urlcolor=azulunir
]{hyperref}

% --- Bibliografía (estilo numérico estándar, sin natbib) ---

% --- Cabecera y pie de página ---
\usepackage{fancyhdr}
\pagestyle{fancy}
\fancyhf{}
\renewcommand{\headrulewidth}{0.4pt}
\fancyhead[L]{\small\textit{Análisis e Interpretación de Datos — UNIR}}
\fancyhead[R]{\small\thepage}

% --- Título compacto ---
\usepackage{titling}
\setlength{\droptitle}{-2em}

% ============================================================
\begin{document}

% ---- PORTADA / CABECERA DEL ARTÍCULO ----------------------
\begin{center}
  {\LARGE\bfseries Impacto del calor extremo sobre la mortalidad en España:\\
  Un análisis estadístico de las temporadas 2024--2025}

  \vspace{0.8em}
  {\large Autor: David Vicente Moya}\\[0.3em]
  {\normalsize Afiliación: Universidad Internacional de La Rioja (UNIR)}\\[0.3em]
  {\normalsize \today}
\end{center}

\vspace{0.5em}
\hrule
\vspace{0.8em}

% ---- RESUMEN ----------------------------------------------
\noindent\textbf{Resumen.}\quad
Este trabajo analiza el impacto del calor extremo sobre la mortalidad en España
durante las temporadas estivales de 2024 y 2025, a partir de datos semanales del
Sistema de Monitorización de la Mortalidad Diaria (MoMo) del Instituto de Salud
Carlos III. Se aplican dos modelos estadísticos:
un análisis descriptivo de la serie temporal, que caracteriza la evolución del exceso
de mortalidad y permite comparar ambas temporadas, y una regresión lineal simple que
cuantifica la relación entre las defunciones atribuibles a la temperatura y el exceso
de mortalidad en las 36 semanas estivales. Los resultados muestran que el verano de
2025 fue un 91\,\% más letal que el de 2024 en términos de defunciones atribuibles
al calor (3.833 frente a 2.008). La regresión lineal obtiene un coeficiente de
determinación $R^2 = 0{,}58$, lo que indica que la temperatura explica el 58\,\%
de la varianza del exceso de mortalidad, con un efecto amplificador estimado de
1{,}07 muertes en exceso por cada defunción directamente atribuida al calor.
Se discuten las limitaciones de ambos enfoques y se proponen líneas futuras basadas
en la incorporación de datos meteorológicos y análisis a escala provincial.
\vspace{0.8em}
\hrule
\vspace{1em}

% ============================================================
\section{Introducción y estado del tema}
% ============================================================

El cambio climático ha intensificado la frecuencia e intensidad de las olas de
calor, convirtiendo el exceso de temperatura en uno
de los principales riesgos medioambientales para la salud pública en España
\cite{inclusion2021}. La evidencia científica ha establecido una relación
entre temperatura y mortalidad: por encima de un umbral térmico
crítico, la mortalidad diaria aumenta de forma aproximadamente lineal con la
temperatura máxima \cite{gasparrini2017}.

En el contexto español, el verano de 2022 supuso un punto de inflexión con más
de 11.000 defunciones atribuidas al calor \cite{momo}. Las temporadas
2024 y 2025 han vuelto a registrar episodios de calor extremo con impacto
significativo en la mortalidad, lo que justifica su estudio sistemático.

El Sistema de Monitorización de la Mortalidad Diaria (MoMo) proporciona
datos semanales de mortalidad observada, mortalidad esperada y exceso
atribuible a temperatura para el conjunto del territorio nacional
\cite{momo}. Esta fuente constituye la base empírica del presente trabajo.

\textbf{Pregunta de investigación:} ¿En qué medida el calor extremo explica el
exceso de mortalidad observado en España durante los veranos de 2024 y 2025,
y qué diferencias se detectan entre ambas temporadas?

% ============================================================
\section{Metodología y resultados obtenidos}
% ============================================================

\subsection{Fuente de datos}

Se han utilizado los datos semanales del sistema MoMo del ISCIII correspondientes
al periodo comprendido entre la semana epidemiológica\footnote{La semana epidemiológica 
es la unidad estándar de vigilancia sanitaria internacional. Se numeran del 1 al 52, 
comienzan el lunes y terminan el domingo. Así, \texttt{2024-31} corresponde a la 
semana del 29~de~julio al 4~de~agosto de 2024.} 20 de 2024 y la semana
39 de 2025 ($n = 72$ semanas). Las variables disponibles son: defunciones
notificadas, defunciones observadas, defunciones estimadas,
exceso de mortalidad sobre todas las causas y defunciones atribuibles a la
temperatura. La base de datos se adjunta como archivo independiente conforme
a los requisitos de la actividad.

Como variable de respuesta se emplea el exceso de mortalidad sobre todas las
causas (diferencia entre defunciones observadas y las estimadas como línea
base) en lugar de las defunciones brutas. Esta elección se justifica porque
el exceso aísla el impacto de las altas temperaturas eliminando la mortalidad
no relacionada, que permanece relativamente constante a lo largo del
año. Usar las defunciones brutas
introduciría ruido que dificultaría la identificación del efecto del calor.

Para el análisis se define la \textit{temporada estival} como las semanas
epidemiológicas 22 a 39 de cada año (aproximadamente de junio a septiembre),
periodo en el que se concentra el riesgo por calor en España.

\subsection{Modelo 1 - Análisis descriptivo de la serie temporal}

El primer modelo consiste en un análisis descriptivo de la serie temporal
completa de 72 semanas, con especial atención a las semanas estivales.
La Figura~\ref{fig:serie} muestra la evolución del exceso de mortalidad
semanal y las defunciones atribuibles al calor.

% --- FIGURA 1: Serie temporal ---
\begin{figure}[H]
\centering
\includegraphics[width=\linewidth]{figura1_serie_temporal.pdf}
\caption{Evolución semanal del exceso de mortalidad y defunciones atribuibles
a la temperatura. España, semanas 20/2024--39/2025 ($n=72$). Fuente: MoMo, ISCIII.}
\label{fig:serie}
\end{figure}

% --- TABLA 1: Estadísticas descriptivas ---
\begin{table}[H]
\centering
\caption{Estadísticos descriptivos por temporada estival (semanas 22--39).}
\label{tab:descriptiva}
\begin{tabular}{lcccc}
\toprule
\textbf{Temporada} & \textbf{$n$ sem.} &
\textbf{Media exceso} & \textbf{DE exceso} &
\textbf{Total atrib. calor} \\
\midrule
Verano 2024 & 18 & 27{,}8  & 292{,}2 & 2.008 \\
Verano 2025 & 18 & 341{,}3 & 320{,}5 & 3.833 \\
\bottomrule
\end{tabular}
\end{table}

El análisis descriptivo revela diferencias sustanciales entre ambas temporadas
(Tabla~\ref{tab:descriptiva}). El verano de 2025 registró una media de exceso
de mortalidad de 341{,}3 defunciones semanales, frente a las 27{,}8 del verano
de 2024, lo que supone un impacto promedio aproximadamente 12 veces superior.

La desviación estándar aporta información adicional sobre la naturaleza del
impacto en cada temporada. En 2024, la DE (292{,}2) supera ampliamente a la
media (27{,}8), lo que refleja un patrón concentrado: semanas con exceso
intenso alternadas con semanas de exceso negativo. En 2025, en cambio, la
DE (320{,}5) es comparable a la media (341{,}3), lo que indica un calor más
sostenido a lo largo de la temporada de verano.

En conjunto, el total acumulado de defunciones atribuibles al calor ascendió
a 3.833 en el verano de 2025 frente a 2.008 en 2024, un incremento del
91\,\% con idéntico número de semanas de observación ($n = 18$).

\subsection{Modelo 2 - Regresión lineal simple}

El segundo modelo aplica una regresión lineal simple sobre las 36 semanas
estivales del periodo de estudio, tomando las defunciones atribuibles a la
temperatura ($X$) como variable independiente y el exceso de mortalidad
sobre todas las causas ($Y$) como variable dependiente.

% --- FIGURA 2: Diagrama de dispersión + recta ---
\begin{figure}[H]
\centering
\includegraphics[width=0.9\linewidth]{figura2_regresion.pdf}
\caption{Diagrama de dispersión y recta de regresión entre defunciones atribuibles a la temperatura y
exceso de mortalidad. Semanas estivales 2024–2025 (n = 36). R2 = 0,58.}
\label{fig:regresion}
\end{figure}

La pendiente estimada ($b_1 = 1{,}074$) indica que, por cada
defunción adicional atribuida directamente a la temperatura, el exceso total
de mortalidad aumenta en aproximadamente 1{,}07 muertes. El valor del
intercepto próximo a cero ($b_0 = 10{,}4$) es coherente con la
expectativa de exceso nulo en ausencia de impacto térmico. El coeficiente
$R^2 = 0{,}58$ señala que el modelo lineal explica el 58\,\% de la varianza
del exceso de mortalidad en semanas estivales.

% ============================================================
\section{Conclusiones y líneas futuras}
% ============================================================

\subsection{Discusión comparada de los modelos}

Ambos modelos apuntan en la misma dirección: el calor extremo tiene un impacto
estadísticamente apreciable sobre la mortalidad en España, con el verano de
2025 sensiblemente más severo que el de 2024. Sin embargo, presentan
discrepancias relevantes que enriquecen el análisis.

El \textbf{modelo descriptivo} revela que varias semanas del verano de 2024
presentan excesos de mortalidad negativos (observadas $<$ esperadas). Esto 
sugiere que las personas más vulnerables fallecen durante el pico de calor,
reduciendo las defunciones en semanas posteriores. Este fenómeno no es captado
por la regresión lineal.

La \textbf{regresión lineal}, con $R^2 = 0{,}58$, confirma una
asociación positiva y sustancial, pero el 42\,\% de varianza inexplicada
apunta a que otros factores modulan el impacto letal de las altas temperaturas. 
Además, la relación entre temperatura y mortalidad no es estrictamente
lineal, sino que sigue una curva, lo que limita la validez
del modelo para valores extremos.

\subsection{Conclusiones principales}

\begin{itemize}
  \item El verano de 2025 registró 3.833 defunciones atribuibles al calor,
        un 91\,\% más que las 2.008 del verano de 2024.
  \item La regresión lineal estima que por cada muerte directamente atribuida
        al calor el exceso total de mortalidad aumenta en 1{,}07 defunciones,
        indicando un efecto amplificador que incluye muertes asociadas pero no
        directamente identificadas como golpe de calor.
  \item Ninguno de los dos modelos por sí solo ofrece una imagen completa:
        el descriptivo aporta contexto temporal y el modelo lineal cuantifica
        la asociación, pero ambos son complementarios y necesarios.
\end{itemize}

\subsection{Limitaciones y líneas futuras}

El principal limitante del análisis realizado es la ausencia de datos de
temperatura en la base de datos utilizada. Incorporar la temperatura
máxima diaria permitiría construir modelos de regresión más precisos 
que son el estándar en epidemiología ambiental
\cite{gasparrini2017}. Como línea futura prioritaria se propone replicar
el análisis a nivel provincial para identificar territorios de mayor
vulnerabilidad, y ampliar la línea temporal a la serie histórica
2015--2025 para evaluar tendencias a largo plazo en el contexto del
cambio climático.

% ============================================================
% BIBLIOGRAFÍA
% ============================================================
\bibliographystyle{plain}
\begin{thebibliography}{9}

\bibitem{inclusion2021}
Ministerio de Inclusi\'{o}n, Seguridad Social y Migraciones (s.f.).
\textit{Cambio clim\'{a}tico: una realidad ya}.
Carta Espa\~{n}a. Gobierno de Espa\~{n}a.
Recuperado el 24 de mayo de 2026, de
\url{https://www.inclusion.gob.es/web/cartaespana/w/cambio-climatico-una-realidad-ya}

\bibitem{gasparrini2017}
Gasparrini, A., \textit{et al.} (2017).
Projections of temperature-related excess mortality under climate change
scenarios. \textit{The Lancet Planetary Health},  2017; 1, e360-e367.
\url{https://doi.org/10.1016/S2542-5196(17)30156-0}

\bibitem{momo}
Instituto de Salud Carlos III.
\textit{Sistema de Monitorización de la Mortalidad Diaria (MoMo).
Panel de datos.}. ISCIII, Madrid.
\url{https://momo.isciii.es}

%\bibitem{romanello2023}
%Romanello, M., \textit{et al.} (2023).
%The 2023 report of the Lancet Countdown on health and climate change.
%\textit{The Lancet}, 402(10419), 2022--2059.
%\url{https://doi.org/10.1016/S0140-6736(23)01859-7}
%
%\bibitem{sanidad2025}
%Ministerio de Sanidad (2025).
%\textit{Plan Nacional de actuaciones preventivas de los efectos
%del exceso de temperaturas sobre la salud. 2025}.
%Direcci\'{o}n General de Salud P\'{u}blica y Equidad en Salud.
%Gobierno de Espa\~{n}a.
%\url{https://www.sanidad.gob.es/areas/sanidadAmbiental/riesgosAmbientales/calorExtremo/publicaciones/docs/balancePlanCalor_2025.pdf}

\end{thebibliography}

\end{document}